% Section 3: The metric tensor and the musical isomorphisms | Part: 1 | Depends on: sec:vectors, sec:tensors | Status: done | Original lines: 778-786
%!TEX root = ../../main.tex
% =====================================================================
\section{The metric tensor and the musical isomorphisms}\label{sec:metric}
% =====================================================================

The dual space $V^*$ exists with no further structure. A \emph{metric} is extra structure: a symmetric non-degenerate
$2$-tensor. It is the only thing that allows us to speak of lengths and angles, and it provides the dictionary
between vectors and covectors (the ``musical isomorphisms'' $\flat$ and $\sharp$). Remember the division of labour
that runs through the document: \emph{covectors, forms and contraction need no metric; gradients, curl, Hodge star and the geometric
product do.}
% Section 3.1: 01-metric-and-musical-isomorphisms | Part: 1 | Depends on: sec:vectors, sec:tensors | Status: done | Original lines: 787-833
%!TEX root = ../../main.tex
% REWRITTEN: notation E -> V, m -> n; the sign conventions for the signature are stated once, in the pitfall of
% Section sec:quadform (the original sentence here pointed to sec:clifford, the pitfall to sec:ga-calculus).
\begin{formal}[{metric}]
A \emph{metric} on $V$ is a tensor $g\in\TT{0}{2}(V)$ (a bilinear map $g:V\times V\to\R$) such that
\begin{enumerate}[label=(\roman*),nosep]
\item $g(v,w)=g(w,v)$ (symmetric);
\item $g(v,w)=0\ \ \forall w\in V\ \Rightarrow\ v=0$ (non-degenerate).
\end{enumerate}
In a basis, $g=g_{ij}\,e^{i}\otimes e^{j}$ with $(g_{ij})$ a symmetric \emph{invertible} matrix. Being symmetric, it can be diagonalized over $\R$; if it has $p$ positive and $q$ negative eigenvalues, the \emph{signature} is $\sig(g)=(p,q)$ (it does not depend on the basis, by Sylvester's law of inertia). 
\begin{itemize}[nosep]
\item $\sig=(n,0)$ (positive definite): \textbf{Riemannian}: $g$ is an inner product.
\item $\sig=(n-1,1)$: \textbf{Lorentzian} (one negative direction, ``mostly plus''), used in GR with $n=4$, $\sig(g)=(3,1)$, \eg Minkowski $\eta=\mathrm{diag}(-1,1,1,1)$. The opposite (``mostly minus'') convention, $\eta=\mathrm{diag}(1,-1,-1,-1)$, and where each one is used, are listed in the pitfall of Section~\ref{sec:quadform}.
\end{itemize}
The \emph{norm squared} is $\|v\|^2=g(v,v)$, which can be positive, zero or negative in the Lorentzian case.
\end{formal}

\begin{intuition}[{the metric is the ``dot product'': it defines lengths, angles and (in GR) causality}]
$g(v,w)$ tells you the length of $v$, the angle between $v$ and $w$, and the projection of $v$ on $w$. In ordinary space all vectors have $g(v,v)\geq0$. In relativity the metric is \emph{indefinite}: a vector $v$ with $g(v,v)<0$ is \textbf{timelike} (a possible velocity of a massive particle), $g(v,v)=0$ with $v\ne0$ is \textbf{null} (light), $g(v,v)>0$ is \textbf{spacelike}. The null vectors form the \emph{light cone} and non-degeneracy does not mean ``$g(v,v)\neq0$ for all $v\neq0$'', it only means that $v\neq0$ always has \emph{some} partner $w$ with $g(v,w)\neq0$.
\end{intuition}

\begin{formal}[{musical isomorphisms (raising and lowering indices)}]
A metric gives an isomorphism between $V$ and $V^*$:
\[
g^\flat:V\to V^*,\quad v\mapsto g(v,\cdot),\qquad\qquad g^\sharp:=(g^\flat)^{-1}:V^*\to V.
\]
In components: $(g^\flat v)_k=v^jg_{jk}=:v_k$ (``lowering the index''), and, with $(g^{kl})$ the inverse matrix of $(g_{ij})$ (so $g^{ij}g_{jk}=\delta^i{}_k$), $(g^\sharp f)^l=f_kg^{kl}=:f^l$ (``raising the index''). The \emph{inverse metric} $g^{-1}=g^{ij}e_i\otimes e_j\in\TT{2}{0}$ is the induced metric on $V^*$. More generally the metric lets us raise or lower any index of any tensor, \eg $T^i{}_j\mapsto T_{ij}=g_{ik}T^k{}_j$; contraction with $g^{ij}$ is the way to take traces of $(0,2)$ tensors.
\end{formal}

\begin{intuition}[{the metric identifies arrows with rulers}]
Given an arrow $v$, the covector $g(v,\cdot)$ is ``the ruler that measures how much of any vector $w$ lies along $v$ (scaled by $|v|$)''. This is the \emph{only} reason why in Euclidean space we can pretend that arrows and covectors are the same thing: with $g=\delta_{ij}$ lowering and raising indices does not change the numbers. With another metric (polar coordinates, curved space, Minkowski) the numbers do change, and the difference between upper and lower indices matters.
\end{intuition}

\begin{example}[{Minkowski space}]
$\eta=\mathrm{diag}(-1,1,1,1)$, $\eta^{-1}=\mathrm{diag}(-1,1,1,1)$ too. For $v=(v^0,v^1,v^2,v^3)$ the lowered components are $v_\mu=(-v^0,v^1,v^2,v^3)$, so
\[
\eta(v,v)=-(v^0)^2+(v^1)^2+(v^2)^2+(v^3)^2.
\]
The vector $(1,1,0,0)$ is null (a light ray with $c=1$), $(2,1,0,0)$ is timelike, $(1,2,0,0)$ is spacelike. For a worldline $x^\mu(\tau)$, \emph{proper time} is $\dd\tau^2=-\eta(\dot x,\dot x)\,\dd\lambda^2$. A Lorentz transformation is precisely a basis change $A$ with $A^t\eta A=\eta$ (compare with $M'=A^tMA$ in Section~\ref{sec:tensors}).
\end{example}

% REORDERED (F2): the example "raising an index in polar coordinates" (original L826-831) moved to Section sec:pushpull.

% Section 3.2: Quadratic forms, polarization and signature | Part: 1 | Depends on: sec:vectors, sec:tensors | Status: done | Original lines: 844-881
%!TEX root = ../../main.tex
\subsection{Quadratic forms, polarization and signature}\label{sec:quadform}
The geometric algebra of Sections~\ref{sec:clifford}--\ref{sec:duality} is built directly on the quadratic form $Q(v)=g(v,v)$
associated with the metric, so we recall the relation between the two and the notation for the signature.
% REWRITTEN: merged with the formal box "metric" above, which is the definition of a metric. Kept here: the
% quadratic form, polarization, the signature (p,q,r) and R^{p,q}. As in the geometric-algebra literature, this
% definition also allows a degenerate g (r>0); the formal box above does not (logged in docs/NOTATION.md).
\begin{definition}[Metric, quadratic form]
Let $g:V\times V\to\K$ be a symmetric bilinear form (a metric as above, but here non-degeneracy is not assumed). It induces the
\emph{quadratic form} $Q(v)=g(v,v)$, and conversely (when $\mathrm{char}\,\K\neq2$) $g$ is
recovered by \emph{polarization}:
\begin{equation}\label{eq:polar}
  g(u,v)=\tfrac12\bigl(Q(u+v)-Q(u)-Q(v)\bigr).
\end{equation}
$g$ is \emph{non-degenerate} if $g(u,\cdot)=0$ implies $u=0$. By Sylvester's law of inertia, over
$\R$ there is an \emph{orthonormal basis} with $g(e_i,e_j)=0$ for $i\neq j$ and
$g(e_i,e_i)\in\{+1,-1,0\}$; the numbers $(p,q,r)$ of $+1,-1,0$ form the \emph{signature}.
We write $\R^{p,q}$ when $r=0$ and $\R^n=\R^{n,0}$ for the Euclidean case.
\end{definition}

% REWRITTEN: the repeated statement of the musical isomorphism (original L861) dropped; see the formal box above.
We shall write $u\cdot v:=g(u,v)$.

\begin{intuition}
The metric is the ruler-and-protractor: it tells you how long vectors are and what angle they make.
Signature $(p,q)$ says that $p$ directions behave like space and $q$ like time (e.g.\ $\R^{1,3}$ for
special relativity). The exterior algebra will need \emph{no} metric; geometric algebra needs one.
\end{intuition}


\begin{pitfall}[three sign conventions for the signature]
(i) Throughout this document the signature $(p,q)$ counts $p$ \emph{positive} and $q$ \emph{negative} eigenvalues,
so that $\R^{p,q}$ and $\Cl(p,q)$ have $p$ basis vectors with $e_i^2=+1$ and $q$ with $e_i^2=-1$; Euclidean space is $(n,0)$.
(ii) General relativity (and Sections~\ref{sec:phys}--\ref{sec:bridge}) uses ``mostly plus'' $\eta=\mathrm{diag}(-1,1,1,1)$, i.e.\ signature
$(3,1)$; (iii) the spacetime algebra of Section~\ref{sec:ga-calculus} uses ``mostly minus'' $\gamma_0^2=+1$, $\gamma_i^2=-1$, i.e.\ $(1,3)$.
Physical statements are the same in the two conventions; only the signs of formulas such as $\vd^2=\pm\Box$ (for the vector
derivative $\vd$ of Section~\ref{sec:ga-calculus}) or
$p^2=\mp m^2$ differ. Always check which one a formula assumes. This box is the reference for signature conventions in the whole document.
\end{pitfall}

% Section 3 (exercises): Exercises | Part: 1 | Depends on: sec:vectors, sec:tensors | Status: done | Original lines: 834-843
%!TEX root = ../main.tex
% REORDERED: moved from the middle of the section to its end. Exercise 3 gets a preview reference (F2).
\begin{exercises}
\begin{enumerate}[nosep]
\item Show that if $g(v,v)>0$ for all $v\neq0$, then $g$ is non-degenerate. Give a symmetric bilinear form on $\R^2$ that is degenerate.
\item In Minkowski space, find $w$ with $\eta(w,w)=0$, $w\ne0$ and $\eta(w,u)\neq0$ for some $u$. Why is this consistent with non-degeneracy?
\item Compute $g^\flat$ of $v=\partial_\varphi$ in polar coordinates (the polar metric $g=\dd r\otimes\dd r+r^2\,\dd\varphi\otimes\dd\varphi$ is derived in Example~\ref{ex:polarmetric}; come back to this exercise after Section~\ref{sec:pushpull}).
\end{enumerate}
\end{exercises}

% Section 3 (keypoints): Keypoints | Part: 1 | Depends on: sec:vectors, sec:tensors | Status: done | Original lines: 882-890
%!TEX root = ../../main.tex
\begin{keypoints}[vectors, tensors and the metric]
\begin{itemize}
\item Vectors have components that transform \emph{against} the basis; covector components transform \emph{with} it; their contraction $\omega_iv^i$ is invariant and metric-free.
\item A tensor of type $(\ell,k)$ is a multilinear machine; $T(V)$ is the free algebra on $V$ and every later algebra is $T(V)$ modulo relations.
\item A metric converts vectors into covectors ($\flat$) and back ($\sharp$); without it there is no ``gradient vector'' and no angle.
\end{itemize}
\end{keypoints}

